
\documentclass[11pt,a4paper]{article}
\usepackage[margin=1in]{geometry}
\usepackage[T1]{fontenc}
\usepackage[utf8]{inputenc}
\usepackage{lmodern}
\usepackage{microtype}
\usepackage{hyperref}
\usepackage{enumitem}
\usepackage{booktabs}
\usepackage{amsmath}
\usepackage{array}
\usepackage{xcolor}
\usepackage{longtable}
\usepackage{tcolorbox}
\usepackage{tikz}
\usetikzlibrary{arrows.meta,positioning,shapes.geometric}
\hypersetup{colorlinks=true, linkcolor=black, urlcolor=blue, citecolor=black}
\setlist[itemize]{topsep=3pt,itemsep=2pt,parsep=0pt}
\setlist[enumerate]{topsep=3pt,itemsep=2pt,parsep=0pt}
\newcommand{\projecttitle}{Semantic Architectural Search for Floor Plans}
\newcommand{\code}[1]{\texttt{\detokenize{#1}}}
\definecolor{darkblue}{RGB}{30,70,120}
\definecolor{lightblue}{RGB}{237,245,252}
\definecolor{lightgraybox}{RGB}{248,248,248}
\tcbset{colback=lightgraybox,colframe=black!20,arc=2mm,boxrule=0.4pt,left=2mm,right=2mm,top=1mm,bottom=1mm}
\title{\textbf{\projecttitle}\\\large Project Report and Implementation Plan}
\author{PrismTeam - CVCS Maticad Project}
\date{July 2026}
\begin{document}
\maketitle
\begin{abstract}
This document defines the complete technical plan for the CVCS Maticad floor plan retrieval project. The immediate objective is apartment-level semantic architectural search: given a complete or partial floor plan query, retrieve the top-k most compatible residential layouts. The plan follows the tutor-provided Phase 1 requirements: start from RPLAN, implement a vision-based retrieval pipeline, implement a graph-based retrieval pipeline inspired by Graph2Plan, test both systems under partial layout degradation, and evaluate retrieval using standard, semantic, and topological metrics. ResPlan is used as a reference for vector-graph data schemas and benchmark design, while HomeWorld is used as a secondary reference for structured JSON representations. The later Maticad stage extends the same retrieval infrastructure to single-room and furniture-layout search.
\end{abstract}
\tableofcontents
\newpage

\section{Executive Summary}

The project builds a \textbf{semantic architectural search system} for residential floor plans. Its first phase focuses on apartment-level retrieval. Given a complete or partial floor plan query, the system retrieves the top-10 most compatible floor plans from a dataset.

The core comparison required by the project is:
\begin{itemize}
    \item \textbf{Vision-based retrieval}: floor plan raster image $\rightarrow$ pretrained image encoder $\rightarrow$ embedding $\rightarrow$ FAISS search.
    \item \textbf{Graph-based retrieval}: floor plan $\rightarrow$ room adjacency graph $\rightarrow$ graph features or GNN embedding $\rightarrow$ FAISS search.
\end{itemize}

The final retrieval system combines both branches through hybrid candidate union and semantic/topological re-ranking.

\begin{tcolorbox}
\textbf{Main Phase 1 objective.} Compare visual and graph representations for retrieving complete, semantically compatible residential floor plans from complete or partial layout queries.
\end{tcolorbox}

\section{Tutor Requirements and Project Priorities}

The tutor document defines the project as \textbf{CVCS Phase 1: Vision vs Graph Representations for Semantic Architectural Search}. The requested direction is operational and specific:
\begin{enumerate}
    \item review and identify relevant floor plan datasets;
    \item start with RPLAN because it provides raster floor plan images;
    \item implement a vision-based retrieval pipeline using CLIP, ViT, DINOv2, DINOv3, or similar pretrained encoders;
    \item extract embeddings and perform nearest-neighbor retrieval with FAISS;
    \item transform floor plan images into structured graphs inspired by Graph2Plan and related graph-based methods;
    \item apply a GNN encoder such as GCN, GAT, GraphSAGE, or similar architectures;
    \item test retrieval on partial floor plans by intentionally degrading complete layouts;
    \item evaluate using Recall@K, semantic room-level similarity, adjacency metrics, functional constraints, and optionally graph edit distance.
\end{enumerate}

The second part of the project concerns \textbf{single-room retrieval} based on furniture layout and spatial arrangement. Public datasets for this exact task are less established, so the proprietary Maticad dataset is expected to become the most relevant source when available.

\section{Research Questions}

The implementation should answer the following questions:
\begin{enumerate}
    \item Does a pretrained visual encoder retrieve floor plans that are visually and geometrically similar?
    \item Does a graph representation retrieve floor plans that are semantically and topologically similar?
    \item Which representation is more robust when the input layout is partial or degraded?
    \item Does a hybrid re-ranking strategy improve retrieval quality over vision-only and graph-only retrieval?
    \item Which graph features and edge definitions best support architectural similarity?
    \item How can the same infrastructure be extended to Maticad single-room and furniture-layout retrieval?
\end{enumerate}

\section{Related Work as Design Input}

\subsection{Graph2Plan}

Graph2Plan is the most important methodological reference for the graph branch. It trains on RPLAN and represents sparse user constraints as layout graphs. Each room is modeled as a graph node, with features for room type, room location, and room size. The room location is discretized using a $5 \times 5$ grid over the building boundary, and room size is encoded through area-related bins or ratios. Edges encode room adjacency and directed spatial relations such as left-of, right-of, above, below, and diagonal variants \cite{graph2plan}.

The original Graph2Plan task is generation: input boundary plus layout graph $\rightarrow$ floor plan raster image plus refined room boxes. For this project, the generation component is not the Phase 1 target. The relevant part is the \textbf{layout graph extraction and retrieval process}:
\[
\text{RPLAN raster/annotations} \rightarrow \text{layout graph} \rightarrow \text{retrieval and comparison}.
\]

Graph2Plan also introduces the idea that sparse constraints can retrieve multiple valid floor plans. This directly supports partial layout retrieval.

\subsection{ResPlan}

ResPlan is a vector-graph dataset of 17,000 residential floor plans. It is useful as a reference for a clean internal schema because it provides both geometry and graph representations. Rooms are nodes with semantic and geometric attributes, while edges encode adjacency types such as door, open-arch, and shared-wall connections \cite{resplan}.

ResPlan also highlights RPLAN's main limitation: although RPLAN is large and convenient for raster vision baselines, its raster format complicates geometric analysis and graph extraction. This reinforces the need for a robust preprocessing module that converts RPLAN samples into a unified JSON and graph representation.

\subsection{HomeWorld}

HomeWorld is relevant as a secondary reference for structured floor plan serialization. It uses a K-D tree JSON representation to serialize room partitions and enable controllable floor plan generation \cite{homeworld}. In this project, K-D tree serialization can be used as an optional structural representation after the required vision and graph retrieval pipelines are implemented.

\section{Datasets}

\subsection{RPLAN}

RPLAN is the primary Phase 1 dataset. It is selected because it provides raster floor plan images and is used by Graph2Plan. It supports both immediate vision retrieval and graph extraction.

RPLAN will be used for:
\begin{itemize}
    \item raster image retrieval;
    \item extraction of room instances and semantic labels;
    \item construction of layout graphs;
    \item generation of partial layout queries;
    \item comparison between vision, graph, and hybrid retrieval.
\end{itemize}

\subsection{ResPlan}

ResPlan should be used as a schema and benchmark reference. If the dataset is available in a usable form, it can become a secondary evaluation dataset because it provides native vector and graph data.

\subsection{Maticad}

The Maticad dataset is expected to be used in the second project stage. The likely task is room-level retrieval based on furniture layout and spatial arrangement.

\subsection{Other Indoor Datasets}

For the Maticad stage, the project should inspect 3D-FRONT, 3D-FUTURE, Structured3D, ZInD, InteriorNet, and CubiCasa5K. The goal is to identify datasets with room-level objects, furniture locations, top-down layouts, and usable spatial relations.

\section{Unified Data Schema}

Every processed sample should be exported into a standard structure:

\begin{verbatim}
data/processed/rplan/<sample_id>/
    image.png
    layout.json
    graph.json
    graph.pt
    preview.png
    masks/
\end{verbatim}

\subsection{Layout JSON}

The internal layout schema should store all geometric and semantic information needed by later modules:

\begin{verbatim}
{
  "id": "000001",
  "source": "RPLAN",
  "image_path": "image.png",
  "width": 256,
  "height": 256,
  "boundary": [[x0, y0], [x1, y1], [x2, y2]],
  "rooms": [
    {
      "id": 0,
      "type": "Bedroom",
      "area": 1234,
      "area_ratio": 0.18,
      "bbox": [x, y, w, h],
      "centroid": [cx, cy],
      "grid_location_5x5": [gx, gy],
      "polygon": [[x0, y0], [x1, y1], [x2, y2]]
    }
  ],
  "edges": [
    {
      "source": 0,
      "target": 1,
      "relation": "adjacent",
      "spatial_relation": "left_of",
      "edge_type": "door_or_proximity"
    }
  ]
}
\end{verbatim}

\subsection{Graph Representation}

The graph representation should be exported both as JSON and as a PyTorch Geometric object:

\begin{verbatim}
Data(
    x=node_feature_tensor,
    edge_index=edge_index_tensor,
    edge_attr=edge_feature_tensor,
    sample_id=sample_id
)
\end{verbatim}

\section{RPLAN Preprocessing}

The first implementation block is dataset preprocessing. The expected pipeline is:

\begin{enumerate}
    \item load RPLAN images and annotations;
    \item inspect annotation channels;
    \item map room labels to a consistent taxonomy;
    \item extract room instance masks;
    \item compute room area, bbox, centroid, and area ratio;
    \item extract boundary mask and boundary polygon if possible;
    \item compute graph edges through doors, shared boundaries, or mask proximity;
    \item save \code{layout.json}, \code{graph.json}, and \code{graph.pt};
    \item generate debug visualizations.
\end{enumerate}

The debug output must include the original raster, room labels, bounding boxes, centroids, and graph overlay.

\section{Graph Construction}

\subsection{Nodes}

Each room is represented as a node. The minimum node attributes are:
\begin{itemize}
    \item normalized room type;
    \item area ratio;
    \item centroid coordinates;
    \item bounding box;
    \item aspect ratio;
    \item $5 \times 5$ grid location.
\end{itemize}

\subsection{Edges}

Edges encode spatial and topological relationships between rooms. The edge extraction module should support:
\begin{itemize}
    \item door-based adjacency when door annotations are available;
    \item shared-boundary adjacency;
    \item proximity-based adjacency;
    \item spatial relation labels.
\end{itemize}

Recommended edge attributes:
\begin{itemize}
    \item \code{source}, \code{target};
    \item \code{edge_type}: door, open-arch, shared-wall, proximity, unknown;
    \item \code{spatial_relation}: left-of, right-of, above, below, diagonal variants, inside, surrounding;
    \item centroid distance;
    \item shared boundary length when available.
\end{itemize}

\section{Vision Retrieval Pipeline}

The first baseline should use pretrained visual embeddings.

\subsection{Pipeline}

\begin{verbatim}
floor plan image
    -> resize and normalize
    -> CLIP / ViT / DINO / ResNet encoder
    -> embedding vector
    -> L2 normalization
    -> FAISS index
    -> top-k retrieval
\end{verbatim}

\subsection{Deliverables}

\begin{itemize}
    \item \code{build_vision_index.py}
    \item \code{retrieve_vision.py}
    \item image embedding matrix
    \item FAISS image index
    \item top-10 visual retrieval grids
    \item vision retrieval metrics
\end{itemize}

\section{Graph Retrieval Pipeline}

The graph branch should be implemented in two stages.

\subsection{Stage 1: Symbolic Graph Feature Retrieval}

Extract a fixed-length graph feature vector:
\begin{itemize}
    \item number of rooms;
    \item room type histogram;
    \item area histogram;
    \item number of edges;
    \item degree distribution;
    \item spatial relation histogram;
    \item edge type histogram;
    \item boundary descriptor if available.
\end{itemize}

Then index these vectors using FAISS.

\subsection{Stage 2: GNN Retrieval}

Train or evaluate a GNN encoder:
\begin{verbatim}
graph.pt
    -> GCN / GraphSAGE / GAT / GIN
    -> global pooling
    -> projection head
    -> graph embedding
    -> FAISS index
\end{verbatim}

The recommended first model is GraphSAGE or GCN because both are easier to control and debug.

\section{Partial Layout Retrieval}

Partial layout retrieval is a central requirement. The system must retrieve complete compatible layouts from degraded or incomplete queries.

\subsection{Raster Masks}

Raster masks should include:
\begin{itemize}
    \item room removal;
    \item rectangular occlusion;
    \item semantic class removal;
    \item structural component removal.
\end{itemize}

\subsection{Graph Masks}

Graph masks should include:
\begin{itemize}
    \item node removal;
    \item edge removal;
    \item room type masking;
    \item room size masking;
    \item room location masking;
    \item subgraph query extraction.
\end{itemize}

\subsection{Training Pairs}

For contrastive training:
\[
\text{partial graph/image of sample } i \quad \leftrightarrow \quad \text{complete graph/image of sample } i
\]

Other floor plans act as negatives.

\section{Hybrid Retrieval and Re-ranking}

The hybrid system merges candidates from image and graph branches:

\begin{verbatim}
top-k vision candidates
+ top-k graph feature candidates
+ top-k GNN candidates
    -> candidate union
    -> score normalization
    -> semantic/topological re-ranking
    -> final top-10
\end{verbatim}

The final score can be defined as:
\[
S = \alpha S_{visual} + \beta S_{graph} + \gamma S_{room} + \delta S_{adj} + \eta S_{area} + \theta S_{boundary}.
\]

Where:
\begin{itemize}
    \item $S_{visual}$ is image embedding similarity;
    \item $S_{graph}$ is graph embedding similarity;
    \item $S_{room}$ measures room type and count compatibility;
    \item $S_{adj}$ measures adjacency compatibility;
    \item $S_{area}$ measures area distribution compatibility;
    \item $S_{boundary}$ measures footprint compatibility.
\end{itemize}

\section{Evaluation Protocol}

\subsection{Retrieval Metrics}

\begin{itemize}
    \item Recall@1, Recall@5, Recall@10;
    \item MRR;
    \item nDCG;
    \item Top-k accuracy.
\end{itemize}

\subsection{Semantic Metrics}

\begin{itemize}
    \item room count accuracy;
    \item room type histogram similarity;
    \item bedroom count match;
    \item bathroom count match;
    \item kitchen/living presence;
    \item area distribution similarity.
\end{itemize}

\subsection{Topological Metrics}

\begin{itemize}
    \item adjacency precision;
    \item adjacency recall;
    \item adjacency F1;
    \item graph edit distance;
    \item degree distribution similarity;
    \item spatial relation histogram similarity;
    \item required adjacency satisfaction.
\end{itemize}

\subsection{Mask Robustness Metrics}

Evaluate each method under mask levels of 10\%, 30\%, 50\%, and 70\%. For each level, report retrieval and compatibility metrics.

\section{Expected Result Tables}

\subsection{Complete Query Retrieval}

\begin{center}
\begin{tabular}{lccccc}
\toprule
Method & R@1 & R@5 & R@10 & MRR & nDCG \\
\midrule
Vision only & -- & -- & -- & -- & -- \\
Graph features & -- & -- & -- & -- & -- \\
GNN graph & -- & -- & -- & -- & -- \\
Hybrid & -- & -- & -- & -- & -- \\
\bottomrule
\end{tabular}
\end{center}

\subsection{Partial Query Retrieval}

\begin{center}
\begin{tabular}{lcccc}
\toprule
Method & Mask 10\% & Mask 30\% & Mask 50\% & Mask 70\% \\
\midrule
Vision only R@10 & -- & -- & -- & -- \\
Graph features R@10 & -- & -- & -- & -- \\
GNN graph R@10 & -- & -- & -- & -- \\
Hybrid R@10 & -- & -- & -- & -- \\
\bottomrule
\end{tabular}
\end{center}

\subsection{Semantic and Topological Compatibility}

\begin{center}
\begin{tabular}{lcccc}
\toprule
Method & RoomSim & Adj-F1 & GED & ConstraintSat \\
\midrule
Vision only & -- & -- & -- & -- \\
Graph features & -- & -- & -- & -- \\
GNN graph & -- & -- & -- & -- \\
Hybrid & -- & -- & -- & -- \\
\bottomrule
\end{tabular}
\end{center}

\section{Maticad Single-Room Retrieval Extension}

The Maticad stage extends the same retrieval logic from apartment-level layouts to room-level furniture arrangements. At that stage:

\begin{itemize}
    \item rooms or furniture objects become the main retrieval units;
    \item object categories, positions, rotations, and spatial relations become graph features;
    \item room top-view images can be indexed with visual encoders;
    \item furniture graphs can be indexed with graph features or GNN embeddings;
    \item hybrid room retrieval can combine image similarity, room geometry, object composition, and furniture spatial relations.
\end{itemize}

A room-level graph can be defined as:
\begin{verbatim}
nodes = furniture objects
edges = object spatial relations
room boundary = global context
\end{verbatim}

\section{Repository Structure}

\begin{verbatim}
PrismTeam/
    data/
        raw/
        processed/
        embeddings/
        indexes/
        splits/
    src/
        data/
        graph/
        vision/
        masking/
        retrieval/
        evaluation/
        visualization/
    scripts/
    notebooks/
    outputs/
    checkpoints/
    docs/
    report/
    README.md
\end{verbatim}

\section{Implementation Roadmap}

\subsection{Phase 0: Repository Setup}

Create the repository structure, environment file, configuration file, and documentation folder.

\subsection{Phase 1: Dataset Audit}

Load RPLAN, inspect images and annotations, map room labels, compute statistics, and create visualization grids.

\subsection{Phase 2: Layout JSON Construction}

Extract room instances, labels, bounding boxes, centroids, area ratios, and boundary information.

\subsection{Phase 3: Layout Graph Extraction}

Build graph JSON and PyTorch Geometric objects. Generate graph overlays for debugging.

\subsection{Phase 4: Vision Retrieval Baseline}

Extract visual embeddings, build a FAISS index, retrieve top-10 floor plans, and evaluate complete image queries.

\subsection{Phase 5: Symbolic Graph Retrieval}

Extract graph feature vectors, build a FAISS graph index, retrieve top-10 layouts, and evaluate semantic/topological similarity.

\subsection{Phase 6: Partial Layout Masking}

Generate raster masks, graph masks, semantic constraint queries, and benchmark splits.

\subsection{Phase 7: Partial Retrieval Evaluation}

Evaluate vision and graph retrieval under different mask types and mask levels.

\subsection{Phase 8: GNN Retrieval}

Train a GNN encoder using full and partial graph pairs, extract embeddings, and evaluate retrieval.

\subsection{Phase 9: Hybrid Retrieval}

Merge candidates, normalize scores, compute semantic and topological scores, and produce final top-10 retrieval results.

\subsection{Phase 10: Final Report and Maticad Preparation}

Prepare final tables, figures, qualitative grids, repository documentation, and the design for the Maticad room-level stage.

\section{First Working Milestones}

\begin{enumerate}
    \item \textbf{Vision demo}: RPLAN image query $\rightarrow$ image encoder $\rightarrow$ FAISS $\rightarrow$ top-10 visually similar floor plans.
    \item \textbf{Graph demo}: RPLAN graph query $\rightarrow$ graph feature vector $\rightarrow$ FAISS $\rightarrow$ top-10 structurally similar floor plans.
    \item \textbf{Partial demo}: masked raster or partial graph query $\rightarrow$ retrieve complete compatible layouts.
    \item \textbf{Hybrid demo}: combine vision and graph results and show final top-10 with interpretable scores.
\end{enumerate}

\section{Final Deliverables}

The project should produce:
\begin{itemize}
    \item RPLAN preprocessing scripts;
    \item internal JSON floor plan representation;
    \item graph extraction module inspired by Graph2Plan;
    \item vision retrieval baseline;
    \item graph feature retrieval baseline;
    \item GNN retrieval model;
    \item partial layout masking benchmark;
    \item hybrid retrieval and re-ranking module;
    \item quantitative evaluation tables;
    \item qualitative top-10 retrieval grids;
    \item GitHub-ready documentation;
    \item final PDF report.
\end{itemize}

\section{Conclusion}

The project should be implemented as a retrieval-first semantic architectural search system. The immediate priority is to build reliable and comparable vision and graph retrieval pipelines on RPLAN, then evaluate their behavior on partial layout queries. Graph2Plan provides the main blueprint for extracting layout graphs from RPLAN, ResPlan provides a clean reference for vector-graph schemas and evaluation tasks, and HomeWorld suggests optional structured JSON extensions. Once the apartment-level retrieval system is stable, the same architecture can be adapted to the Maticad single-room furniture-layout retrieval task.

\begin{thebibliography}{9}
\bibitem{cvcs} CVCS Maticad. \textit{CVCS Phase 1: Vision vs Graph Representations for Semantic Architectural Search}. Tutor-provided project document, 2026.
\bibitem{graph2plan} R. Hu, Z. Huang, Y. Tang, O. van Kaick, H. Zhang, and H. Huang. \textit{Graph2Plan: Learning Floorplan Generation from Layout Graphs}. ACM Transactions on Graphics, 2020.
\bibitem{resplan} M. Abouagour and E. Garyfallidis. \textit{ResPlan: A Large-Scale Vector-Graph Dataset of 17,000 Residential Floor Plans}. arXiv:2508.14006, 2025.
\bibitem{homeworld} W. Li et al. \textit{HomeWorld: A Unified Floorplan-to-Furnished Framework for Generating Controllable, Densely Interactive Whole-Home Scenes}. arXiv, 2026.
\bibitem{rplan} W. Wu, X.-M. Fu, R. Tang, Y. Wang, Y.-H. Qi, and L. Liu. \textit{Data-driven Interior Plan Generation for Residential Buildings}. ACM Transactions on Graphics, 2019.
\end{thebibliography}

\end{document}
